\section{Evaluation}
\label{sec:evaluation}

To ensure the reproducibility of our empirical results, the source code, devicetree overlays, and configuration files used for the concurrent RTOS workload experiments in this section are publicly available as an open-source artifact at \url{https://github.com/heronet/BASE-IoT-Evaluation}.

\subsection{Memory Footprint}

As shown in Table~\ref{tab:implementation_metrics}, BASE requires a minimum of \SI{3.6}{\kilo\byte} Flash and \SI{644}{\byte} RAM for a complete fingerprint subsystem and driver (ZFM-X0).

\subsection{Abstraction Overhead}

Static inline wrappers add only \numrange{5}{10} CPU cycles per call. Sensor processing time (hundreds of milliseconds) dominates, making the overhead negligible in practice.

\subsection{Portability Analysis}

To quantify portability, we measured application code changes required for two sensor migrations: within-modality (fingerprint to fingerprint) and cross-modality (fingerprint to face/palm). Both used the legacy Sensor API as a baseline. Results are in Table~\ref{tab:sensor_migration}.

\begin{table}[htbp]
\renewcommand{\arraystretch}{1.4}
\caption{Sensor Migration Code Change Analysis}
\label{tab:sensor_migration}
\begin{tabularx}{\linewidth}{|l|X|X|X|}
\hline
\textbf{Change} & \textbf{Legacy API} & \multicolumn{2}{c|}{\textbf{BASE}} \\
\cline{3-4}
 & \textbf{R502A$\to$GT-521F32} & \textbf{R502A$\to$GT-521F32} & \textbf{GT-521F32$\to$AI10} \\
\hline
Headers & Vendor replace & No change & No change \\
\hline
Enrollment & Rewrite for 3 samples & Auto-adapts & Auto-adapts \\
\hline
Attributes & 8 macros replace & No change & No change \\
\hline
Operations & Hardware fns replace & Unified API & Unified API \\
\hline
Devicetree & Update string & Update string & Update string \\
\hline
\textbf{App LOC} & \textbf{47 changed} & \textbf{0 changed} & \textbf{0 changed} \\
\hline
\end{tabularx}
\end{table}

The GT-521F32 to AI10 migration crosses a modality boundary: from a contact fingerprint sensor to a camera-based device reporting two modalities. The application enrollment and match loop required no changes. The \texttt{samples\_required} field adapts from 2 (ZFM-X0) to 3 (GT-5X) to 1 (AI10) transparently. The \texttt{biometric\_match\_result.modality} field allows the application to distinguish face from palm results on the AI10 if needed, without any conditional compilation or sensor-specific headers.

\subsection{Concurrent RTOS Workload Analysis}
\label{sec:eval_concurrency}

To empirically validate BASE's schedulability claims, we measured deadline adherence of a high-priority control task ($\tau_1$) while a low-priority biometric task ($\tau_2$) ran a \SI{500}{\milli\second} capture sequence on the STM32F446RE.

\textbf{Setup:} $\tau_1$ (priority 2) executes a \SI{10}{\milli\second} periodic loop and records its actual period using \texttt{k\_uptime\_ticks()}. A deadline miss is counted when the measured period exceeds the \SI{10}{\milli\second} target by more than \SI{1}{\milli\second}. $\tau_2$ (priority 8) runs concurrently in two configurations:

\begin{itemize}
    \item \textbf{Legacy baseline:} \texttt{k\_sched\_lock()} followed by \\ \texttt{k\_busy\_wait(500000)}, reproducing the behavior of a vendor polling library that disables preemption while spinning on a UART receive loop.
    \item \textbf{BASE:} \texttt{biometric\_enroll\_capture()} with a \SI{600}{\milli\second} timeout, which yields the CPU to the scheduler throughout the capture window via cooperative \texttt{k\_msleep} polling and semaphore-based UART synchronization.
\end{itemize}

\textbf{Results:} Table~\ref{tab:concurrency} summarizes the measurements over 200 samples.

\begin{table}[!ht]
\caption{Concurrent Workload: $\tau_1$ Deadline Adherence (200 samples)}
\label{tab:concurrency}
\centering
\renewcommand{\arraystretch}{1.3}
\begin{tabularx}{\columnwidth}{|l|X|X|X|}
\hline
\textbf{Mode} & \textbf{Max Jitter} & \textbf{Avg Jitter} &
\textbf{Missed} \\
\hline
Legacy (busy-wait) & \SI{490}{\milli\second} & \SI{167}{\milli\second}
& 68 / 200 \\
\hline
BASE (yielding) & \SI{100}{\micro\second} &
\SI{100}{\micro\second} & 0 / 200 \\
\hline
\end{tabularx}
\end{table}

Under the legacy baseline, $\tau_2$ monopolized the CPU for the full \SI{500}{\milli\second} capture window on each iteration, causing $\tau_1$ to miss 68 of 200 deadlines with a maximum jitter of \SI{490}{\milli\second}. Under BASE, $\tau_2$ yielded the CPU on each \texttt{enroll\_capture} call via cooperative polling and semaphore-based UART synchronization. $\tau_1$ missed zero deadlines across all 200 samples, with a maximum jitter of \SI{100}{\micro\second}, consistent with a single scheduling tick of overhead at the configured tick rate. This confirms that BASE allows heavy biometric processing to coexist with hard real-time tasks without deadline violation.

\subsection{Multi-Modal Validation}
\label{sec:eval_multimodal}

To validate the modality independence claim, the enrollment and match application was executed against all three driver families without modifying application code. The DFRobot AI10 driver was developed independently by Benjamin Cabé. We reviewed the driver source and confirmed that zero modifications to the subsystem layer or API headers were required. Physical validation was performed on his hardware, with palm vein recognition demonstrated publicly in the official Zephyr 4.4 release \cite{zephyr_44_release}.

\begin{small}
\begin{verbatim}
biometric_get_capabilities(dev, &caps);
biometric_enroll_start(dev, id);
do {
  biometric_enroll_capture(dev, K_SECONDS(5), &res);
} while (res.samples_captured < res.samples_required);
biometric_enroll_finalize(dev);
biometric_match(dev, BIOMETRIC_MATCH_IDENTIFY,
                0, K_SECONDS(10), &match);
\end{verbatim}
\end{small}

This loop ran without modification against a ZFM-60 (optical fingerprint), a GT-511C3 (optical fingerprint), and a DFRobot AI10 (face and palm). Table~\ref{tab:multimodal_validation} summarizes what was validated on each device.

\begin{table}
\caption{Multi-Modal Validation Summary}
\label{tab:multimodal_validation}
\centering
\renewcommand{\arraystretch}{1.3}
\setlength{\tabcolsep}{3pt}
\begin{tabularx}{\columnwidth}{|X|c|c|c|}
\hline
\textbf{Operation} & \textbf{ZFM-X0} & \textbf{GT-5X} & \textbf{AI10} \\
\hline
\texttt{get\_capabilities} & \checkmark & \checkmark & \checkmark \\
\hline
\texttt{enroll\_start/capture/finalize} & \checkmark & \checkmark & \checkmark \\
\hline
\texttt{match} (identify mode) & \checkmark & \checkmark & \checkmark \\
\hline
\texttt{template\_delete} & \checkmark & \checkmark & \checkmark \\
\hline
\texttt{supported\_modalities} bitmask & \checkmark & \checkmark & \checkmark \\
\hline
\texttt{result.modality} field & N/A & N/A & \checkmark \\
\hline
Timeout/error recovery & \checkmark & \checkmark & \checkmark \\
\hline
\textbf{App LOC changed} & \multicolumn{3}{c|}{\textbf{0}} \\
\hline
\end{tabularx}
\end{table}

The \texttt{samples\_required} field returned 2 for ZFM-X0, 3 for GT-5X, and 1 for the AI10, with the application loop handling all three without any conditional logic. On the AI10, \texttt{result.modality} correctly reported \texttt{BIOMETRIC\_MODALITY\_FACE} for face enrollments and \texttt{BIOMETRIC\_MODALITY\_PALM} for palm enrollments.

\subsection{Limitations and Scope}

While we addressed the core integration problems like vendor fragmentation, type safety violations, and broken state management on resource-constrained RTOSes with BASE, several related problems remain outside the current scope:

\textbf{Template format abstraction}: Templates remain sensor-specific binary data. BASE does not abstract or normalize the template format itself, only the workflow for acquiring and storing them. Cross-sensor template migration is not supported. On sensors that support it, \texttt{biometric\_enroll\_from\_image()} enables enrollment from pre-captured image data, supporting remote provisioning workflows where live capture is impractical.